\respcap{cap:emp}
\resp{em-def} $\{23,45\}$ é maximal, pois toda aresta do caminho tem um extremo em $\{2,3,4,5\}$, mas tem só 2 arestas. $\{12,34,56\}$ é máximo e perfeito: satura os 6 vértices. Com 7 vértices não há emparelhamento perfeito, porque o número de vértices é ímpar; o máximo tem 3 arestas, por exemplo $\{12,34,56\}$, e deixa o vértice 7 livre.

\resp{em-metade} Seja $M^*$ um emparelhamento máximo. Como $M$ é maximal, toda aresta de $M^*$ tem pelo menos um extremo saturado por $M$; e cada aresta de $M$ satura só 2 vértices, cada um deles extremo de no máximo uma aresta de $M^*$. Logo $|M^*|\le2|M|$. A igualdade vale no caminho $1,2,3,4$ com $M=\{23\}$, que é maximal, enquanto $\nu=2$ (com $\{12,34\}$).

\resp{em-perfeitos} Cada vértice é extremo de exatamente uma aresta de $M$ e de uma de $M'$. Se essas duas arestas são a mesma, ela sai da diferença simétrica e o vértice fica com grau 0; senão, ele tem grau 2 em $M\triangle M'$, com uma aresta de cada emparelhamento. As componentes não triviais são, portanto, ciclos (não há extremos de grau 1) com arestas alternadamente em $M$ e $M'$, e um ciclo alternado tem comprimento par. Como $M\ne M'$, há pelo menos um ciclo. Uma árvore não tem ciclos; logo, se tivesse dois emparelhamentos perfeitos distintos, $M\triangle M'$ seria vazio, isto é, $M=M'$, uma contradição.

\resp{em-aumenta} Os vértices livres são 1 e 6, e $P=1,2,3,4,5,6$ alterna $12\notin M$, $23\in M$, $34\notin M$, $45\in M$ e $56\notin M$: é aumentante. $M\triangle P=\{12,34,56\}$, que satura os 6 vértices: é perfeito e, portanto, máximo (nenhum emparelhamento tem mais que $n/2$ arestas).

\resp{em-hung} As buscas a partir de $a$, $b$, $c$ e $d$ encontram vértices livres imediatamente: $a,1$; $b,2$; $c,3$; $d,4$. A busca a partir de $e$ vai a 1, saturado, segue para $a$, depois para 2 e $b$, para 3 e $c$, para 4 e $d$, e encontra 5, livre: o caminho aumentante é $e,1,a,2,b,3,c,4,d,5$, com 9 arestas. O emparelhamento final é $\{e1,a2,b3,c4,d5\}$, perfeito.

\resp{em-hungtree} Na árvore húngara $T$, todo vértice de $Y$ é saturado e o seu par está em $T$, e todo vértice de $X$ em $T$ tem todos os vizinhos em $T$ (senão a busca os teria alcançado); o único vértice livre de $T$ é $r$. Num grafo bipartido, ao longo de um caminho alternante que começa num vértice livre de $X$, os vértices de $Y$ são alcançados por arestas fora de $M$ e os de $X$, por arestas de $M$. Seja $P$ um caminho aumentante que toca $T$. Ele não pode entrar em $T$ por um vértice de $X$ (todos os vizinhos já estão em $T$); se entra por um $y\in Y$, a aresta seguinte é a de $M$, para o par de $y$, que está em $T$, e daí em diante $P$ não sai mais de $T$. Então $P$ teria de terminar num vértice livre de $T$, que só pode ser $r$; mas, lido a partir de $r$, esse caminho chegaria a $y$ por uma aresta de $M$, o que a paridade proíbe. Logo nenhum caminho aumentante usa vértices de $T$. Uma ampliação por um caminho que evita $T$ não muda as arestas de $M$ incidentes em $T$ nem as arestas entre vértices de $T$, e o argumento continua valendo; por isso $T$ pode ser removida.

\resp{em-hall} Num grafo bipartido $k$-regular, $|X|=|Y|$, porque o número de arestas é $k|X|=k|Y|$. Dado $S\subseteq X$, as $k|S|$ arestas que saem de $S$ chegam a $N(S)$, cujos vértices recebem, no total, $k|N(S)|$ arestas; logo $k|S|\le k|N(S)|$ e $|N(S)|\ge|S|$. Pelo teorema de Hall há um emparelhamento que satura $X$, e, como $|X|=|Y|$, ele é perfeito.

\resp{em-hallex} Não. O conjunto $S=\{a,b,c\}$ tem $N(S)=\{1,2\}$, e $|N(S)|=2<3$. O máximo tem 3 arestas, por exemplo $\{a1,b2,d3\}$: um dos três vértices $a,b,c$ sempre fica livre.

\resp{em-konig} Com $M=\{a1,b2,d3\}$, o único vértice livre de $X$ é $c$. As buscas alternantes a partir de $c$ alcançam 1 e 2 (por arestas fora de $M$) e os seus pares, $a$ e $b$: $Z=\{c,1,2,a,b\}$. A cobertura é $C=(X\setminus Z)\cup(Y\cap Z)=\{d\}\cup\{1,2\}=\{d,1,2\}$, com $3=\nu$ vértices: as arestas de $a$, $b$ e $c$ chegam a 1 ou 2, e as de $d$ tocam $d$.

\resp{em-domino} Pinte o tabuleiro como de costume. Cada dominó cobre uma casa branca e uma preta, e os dois cantos opostos têm a mesma cor; sobram 32 casas de uma cor e 30 da outra, e 31 dominós cobririam 31 de cada. No grafo bipartido cujos vértices são as casas e cujas arestas ligam casas vizinhas (os lados são as duas cores), uma cobertura por dominós é um emparelhamento perfeito, que exige lados de mesmo tamanho.

\resp{em-arvore} Seja $f$ uma folha e $p$ o seu vizinho, e $M^*$ um emparelhamento máximo. Se $fp\notin M^*$, então $p$ é saturado por alguma aresta $pq$ (senão $M^*+fp$ seria maior), e $M^*-pq+fp$ é um emparelhamento do mesmo tamanho que contém $fp$. Algoritmo: enquanto houver aresta, escolha uma folha $f$, ponha $fp$ em $M$ e remova $f$ e $p$ (e os vértices que ficarem isolados); com uma fila de folhas e os graus atualizados, custa $O(n)$. Na árvore dada: a folha 4 dá a aresta $24$ (e o 5 fica isolado); a folha 7 dá $67$; o 3 vira folha e dá $13$. O máximo tem 3 arestas, $\{24,67,13\}$.

\resp{em-conta} Em $K_{n,n}$, um emparelhamento perfeito é uma bijeção entre os lados: $n!$. Em $K_{2n}$, o primeiro vértice escolhe o seu par entre $2n-1$, o próximo vértice livre entre $2n-3$, e assim por diante: $(2n-1)(2n-3)\cdots1=\frac{(2n)!}{2^n\,n!}$. Para $n=3$: $3!=6$ e $5\cdot3\cdot1=15$.

\resp{em-ciclo} Nos dois casos, $\nu=\floor{n/2}$: tome arestas alternadas ao longo do caminho ou do ciclo; não dá para fazer melhor, pois cada aresta usa 2 dos $n$ vértices. Há emparelhamento perfeito exatamente quando $n$ é par (para o ciclo, $n\ge3$ par; para o caminho, $n\ge2$ par).

\resp{em-cobertura} Dado um emparelhamento máximo $M$, os $n-2|M|$ vértices livres formam um conjunto independente (senão $M$ cresceria), e cada um tem uma aresta para um vértice saturado; $M$ mais uma aresta por vértice livre cobre tudo, com $|M|+(n-2|M|)=n-\nu$ arestas. Reciprocamente, numa cobertura mínima $C$ cada componente é uma estrela (uma aresta cujos dois extremos têm outra aresta de $C$ seria desnecessária); com $c$ componentes, $|C|=n-c$, e uma aresta de cada estrela forma um emparelhamento de tamanho $c=n-|C|$, logo $\nu\ge n-|C|$. Caminho com 6 vértices: $\nu=3$, cobertura mínima com 3 arestas, $\{12,34,56\}$. Estrela com 5 vértices: $\nu=1$, cobertura mínima com as 4 arestas.

\resp{em-triangulo} $\nu\le\tau$: as arestas de um emparelhamento não têm extremos comuns, e cada uma precisa de um vértice diferente da cobertura. $\tau\le2\nu$: os $2\nu$ extremos de um emparelhamento máximo formam uma cobertura, pois uma aresta sem extremo nesse conjunto poderia ser acrescentada. No triângulo, $\nu=1$ (duas arestas sempre se tocam) e $\tau=2$ (um vértice só não toca a aresta oposta): $\nu<\tau$.

\resp{em-tutte} Seja $M$ um emparelhamento. Em cada componente ímpar de $G-S$, algum vértice fica livre ou é emparelhado com um vértice de $S$, pois os vértices da componente não podem ser todos emparelhados entre si. Como cada vértice de $S$ atende no máximo uma componente, pelo menos $\mathrm{imp}(G-S)-|S|$ vértices ficam livres, e $2|M|\le n-\mathrm{imp}(G-S)+|S|$. Na estrela com 4 vértices, com $S$ igual ao centro, $G-S$ tem 3 componentes ímpares (as folhas): $\nu\le\frac12(4+1-3)=1<2$, e não há emparelhamento perfeito.

\resp{em-flor} $r$ é par. De $r$: $a$ é saturado e sem rótulo, logo $a$ é ímpar e $b$ é par; também $d$ é ímpar e $c$ é par. De $b$: o vizinho $c$ é par, e a aresta $bc$ fecha o ciclo ímpar $r,a,b,c,d$, cuja base é a própria raiz $r$ (a haste é vazia). Contraído o blossom num vértice $B$, livre, a aresta $Bx$ (que vinha de $c$) leva ao vértice livre $x$: o caminho aumentante em $G/B$ é $B,x$. Na expansão, a base $r$ é ligada a $c$ pelo lado par do ciclo, $r,d,c$ ($rd\notin M$, $dc\in M$), e o caminho em $G$ é $r,d,c,x$. O novo emparelhamento é $\{ab,rd,cx\}$, perfeito.

\resp{em-bipcorreto} Com raiz $r\in X$, um caminho de $r$ até um vértice de $X$ tem comprimento par, e até um de $Y$, ímpar, porque cada aresta troca de lado. Na árvore alternante, os vértices pares (a raiz e os alcançados por arestas de $M$) estão a uma distância par de $r$ na árvore e, portanto, estão todos em $X$. Uma aresta entre dois vértices pares ligaria dois vértices de $X$, o que não existe num grafo bipartido. Sem essa aresta não há ciclo ímpar fechado na árvore, nenhum blossom aparece, e cada vértice tem uma única paridade possível: rotulá-lo uma vez basta.

\resp{em-atrib} Os custos das 6 designações são $5,6,6,7,9,11$. O mínimo é \textbf{5}, com $T_1\to2$, $T_2\to1$ e $T_3\to3$ (custos $1+2+2$). Com 20 técnicos seriam $20!\approx2{,}4\cdot10^{18}$ designações. O algoritmo de Kuhn e Munkres (método húngaro com pesos) encontra a designação ótima em $O(n^3)$.

\resp{em-modela} Ponha os alunos de um lado e, do outro, $c_j$ cópias do projeto $j$; ligue cada aluno a todas as cópias dos projetos que ele aceita. Um emparelhamento é uma alocação que respeita as vagas (cada cópia é uma vaga), e um emparelhamento máximo aloca o maior número possível de alunos. Com $a$ alunos, $C=\sum_jc_j$ vagas e $m'$ arestas no grafo com cópias (cada indicação de projeto $j$ vira $c_j$ arestas), o método húngaro faz no máximo $a$ buscas de custo $O(a+C+m')$: $O(a\,(a+C+m'))$.
